\documentclass{article}

\usepackage{iclr2027_conference,times}
\input{math_commands.tex}
\usepackage{booktabs}
\usepackage{hyperref}
\usepackage{url}

\title{Minimal Recovery-Set Topology of Jailbreak Prompts}

% Keep the submission anonymous. Do not insert author identities here.
\author{Anonymous authors}

% Uncomment only for the camera-ready version.
% \iclrfinalcopy

\newcommand{\resultpending}{\textbf{[RESULT PENDING]}}

\begin{document}

\maketitle

\begin{abstract}
Jailbreak prompts are often treated as indivisible strings or ranked by
single-token importance, even when their effects may depend on interactions
among multiple human-editable components. We define a behavioral intervention
problem that recovers every minimal set of attack-added input units whose
neutralization restores safe behavior while preserving the underlying request.
Our protocol requires strict-subset minimality, stability across neutralizers
and decoding seeds, payload invariance, capability controls, and blinded human
audit. \resultpending{} This draft does not claim an empirical finding until the
fresh confirmatory experiment and audit are complete.
\end{abstract}

\section{Introduction}

Understanding which parts of a successful jailbreak are behaviorally necessary
is useful for evaluation, diagnosis, and defense. Existing token scores and
leave-one-out analyses can miss interactions: two units may be jointly
necessary, or several distinct minimal pathways may independently sustain an
attack. We therefore study the topology of robust minimal recovery sets under
direct prompt interventions.

The intended contributions are:
\begin{itemize}
  \item a precise, intent-preserving definition of robust minimal recovery;
  \item an exact coarse-unit enumeration protocol with neutralizer, seed, and
        capability controls; and
  \item \resultpending{} a repeated empirical finding on a fresh, multi-model
        confirmation split.
\end{itemize}

\section{Problem Formulation}

Let $p$ be a harmful payload, $a(p)$ an attacked prompt that contains $p$
exactly once, and $U=\{u_1,\ldots,u_m\}$ the frozen attack-added units. For a
subset $S\subseteq U$ and neutralizer $n$, let $I_{S,n}(a(p))$ replace only the
units in $S$ while leaving $p$ byte-identical. A recovery set must restore the
predeclared safe-behavior predicate across required seeds and neutralizers,
without truncation, malformed input, or general capability collapse. It is
minimal if none of its strict subsets is a recovery set.

The collection of all minimal recovery sets forms a hypergraph over $U$. This
representation exposes non-singleton dependencies and multiple alternative
pathways that a single importance ranking cannot express.

\section{Experimental Design}

\subsection{Eligibility and stable pairs}

Direct and attacked prompts share payload, decoding configuration, and seeds.
A pair is eligible only when the direct request remains non-harmful, the attack
produces substantive harmful assistance, the payload occurs once and is
byte-identical, and neither truncation nor capability failure determines the
label.

\subsection{Interventions}

We freeze at most six coarse attack-added units per eligible pair and enumerate
all $2^m$ subsets. Each subset is evaluated under two independently meaningful
neutralizers and three decoding seeds. Decision-critical outputs enter blinded
human audit by two independent annotators, with a third adjudicator for
disagreement.

\subsection{Baselines and controls}

We compare exact topology with singleton leave-one-out attribution, random
cost-matched subsets, and greedy backward elimination. Payload invariance,
grammar/token validity, matched irrelevant edits, and benign capability probes
separate genuine safety recovery from destructive editing.

\subsection{Models, attacks, and data}

\resultpending{} Model revisions, quantization status, attack families, fresh
payload counts, exclusions, and all artifact hashes will be inserted from the
frozen confirmatory contract. Development examples are excluded from primary
paper-valid estimates.

\section{Results}

\resultpending{} This section will report stable-pair yield, minimal-set size,
multiple-pathway frequency, singleton miss rate, neutralizer agreement, seed
stability, capability-confounded recovery, abstention, and paired uncertainty.
No value will be inserted without a config-linked result artifact and completed
human audit.

\begin{table}[t]
  \caption{Primary confirmatory outcomes. All cells are intentionally pending.}
  \label{tab:primary-results}
  \centering
  \begin{tabular}{lc}
    \toprule
    Outcome & Estimate \\
    \midrule
    Eligible stable attack instances & \resultpending \\
    Instances with robust minimal recovery & \resultpending \\
    Non-singleton or multiple-pathway instances & \resultpending \\
    Singleton attribution miss rate & \resultpending \\
    Cross-neutralizer agreement & \resultpending \\
    Capability-confounded recovery rate & \resultpending \\
    \bottomrule
  \end{tabular}
\end{table}

\section{Related Work}

\resultpending{} Add verified citations on jailbreak evaluation, prompt-level
causal localization, adversarial suffix optimization, and safety classifiers.
The final text will distinguish behavioral recovery topology from internal
representation localization and ordinary token saliency.

\section{Limitations and Ethics}

The intervention lattice is defined over a frozen human-editable partition and
does not establish a unique natural-language causal decomposition. Results may
depend on model revision, quantization, chat template, neutralizer, and decoding
distribution. Behavioral recovery is not evidence that the model has forgotten
the harmful capability. The study minimizes dissemination risk by keeping raw
harmful payloads and operational outputs out of the public repository and by
releasing only hashes, aggregates, and appropriately redacted examples.

\section{Reproducibility Statement}

Every reported result will be linked to a frozen configuration, git commit,
model and tokenizer revision, generation parameters, seeds, hardware record,
and exclusion ledger. The repository includes CPU-only tests and harmless smoke
paths; restricted artifacts remain private and are identified by digest.

\section{AI Use Statement}

Generative AI tools assisted with repository inspection, planning, code and
test maintenance, and drafting of protocol and manuscript text. The human
authors are responsible for the research design, execution, verification,
interpretation, citations, and final submitted content. This statement will be
updated to enumerate material AI use in the final paper.

\section{Conclusion}

We introduced a protocol for measuring robust minimal recovery-set topology in
jailbreak prompts. \resultpending{} A conclusion about empirical topology will
be written only after the fresh confirmatory split and blinded audit satisfy the
predeclared decision rule.

% Enable the bibliography after verified references have been added.
% \bibliographystyle{iclr2027_conference}
% \bibliography{references}

\appendix

\section{Registered Decision Rules}

The appendix will reproduce the final eligibility, GO/NARROW/STOP, exclusion,
and audit rules and provide artifact identifiers. Appendix material will not be
used to hide evidence required for the main claims.

\end{document}

